\documentclass[11pt]{article}
\usepackage[margin=1in]{geometry}
\usepackage{amsmath,amssymb}
\usepackage{booktabs}
\usepackage{array}
\usepackage[hidelinks]{hyperref}
\usepackage{parskip}
\emergencystretch=2em

\newcommand{\B}{\mathcal{B}}
\newcommand{\Leff}{\Lambda_{\rm eff}}
\newcommand{\Ctrue}{\mathcal{C}_{\rm true}}
\newcommand{\Ccg}{\mathcal{C}_{\rm cg}}
\newcommand{\Ciso}{\mathcal{C}_{\rm iso}}
\newcommand{\Cbin}{\mathcal{C}_{\rm iso}^{\rm bin}}
\newcommand{\Tcg}{T_{\rm cg}}

\title{Radiative Cooling on Coarse-Grained TRML Fields:\\
$\Ccg$, $\B[n^2\Lambda(T)]$ and the Isobaric PDF Estimator $\Ciso$}
\date{}

\begin{document}
\maketitle

\section{Setup and notation}

The fine simulation is a 2D turbulent radiative mixing layer (TRML) at $512\times1024$ on a $10\times20$~pc box, following the setup of Sharma et al.\ (2025) but in 2D. The fine cell size is $\delta x = 10\,{\rm pc}/512 = 0.0195$~pc. Coarse-graining by a factor $\Delta$ (in cells) gives coarse cells of physical size $\ell=\Delta\,\delta x$; $\Delta=32$ gives a $16\times32$ grid with $\ell=0.625$~pc, and $\Delta=64$ gives $8\times16$ with $\ell=1.25$~pc. In 2D, ``volume'' means area throughout.

\paragraph{Block mean.} For a field $f$,
\begin{equation}
\B[f]_{ij} \equiv \mathrm{BlockMean}(f,\Delta) = \frac{1}{\Delta^2}\sum_{n=0}^{\Delta-1}\sum_{m=0}^{\Delta-1} f(i\Delta+n,\,j\Delta+m).
\end{equation}
$\B$ is linear and conserves integrals: $\sum_{ij}\B[f]_{ij}\,\ell^2=\int f\,dA$.

\paragraph{Coarse fields.} Conserved and extensive quantities are volume averaged; specific quantities are mass (Favre) averaged:
\begin{align}
\bar\rho &= \B[\rho], &
\bar u_i &= \B[\rho u_i]/\bar\rho, &
\bar P &= \B[P], &
\bar s &= \B[\rho s]/\bar\rho .
\end{align}
The internal energy density $P/(\gamma-1)$ is extensive, so volume averaging it is correct, but it is not separately conserved. The coarse state therefore satisfies $\B[E]=\bar E_{\rm int}+\tfrac12\bar\rho|\bar u|^2+k_{\rm sgs}$, with subgrid kinetic energy $k_{\rm sgs}=\tfrac12(\B[\rho|u|^2]-\bar\rho|\bar u|^2)\ge0$. Averaging $P$ directly keeps $k_{\rm sgs}$ out of the temperature; reconstructing $P$ from $\B[E]$ would raise $\Tcg$ by $(\gamma-1)\mu m_p k_{\rm sgs}/(k_B\bar\rho)$.

\paragraph{Coarse temperature.} From the equation of state applied to coarse fields,
\begin{equation}
\Tcg = \frac{\mu m_p\,\bar P}{k_B\,\bar\rho} = \frac{\B[\rho T]}{\B[\rho]},
\end{equation}
i.e.\ $\Tcg$ is the Favre (mass-weighted) temperature, biased toward the cold phase. By contrast $\B[T]$ is volume weighted and biased toward the hot phase. If the block is isobaric, $\Tcg$ is exactly the harmonic volume mean,
\begin{equation}
\Tcg^{-1} = \B[T^{-1}] \qquad (P=\text{const in block}).
\label{eq:harmonic}
\end{equation}

\paragraph{Cooling function.} As in the reference setup, cooling is switched off in both stable phases:
\begin{equation}
\Leff(T)=\Lambda(T)\,\Theta(T-1.05\,T_c)\,\Theta(0.95\,T_h-T).
\end{equation}
For $T_c=10^4$~K and $T_h=10^6$~K ($\chi=T_h/T_c=100$) the window is $4.021<\log_{10}T<5.978$. Throughout, $n=\rho/(\mu m_p)$ with constant $\mu$, and the kernel must use exactly the simulation's $\Lambda$ table, interpolation, density convention ($n^2$ vs.\ $n_en_H$) and cutoffs.

\section{The three cooling estimates}

\begin{align}
\Ctrue &\equiv \B\!\left[n^2\Leff(T)\right] && \text{true block-averaged cooling,}\\
\Ccg   &\equiv \bar n^2\,\Leff(\Tcg),\quad \bar n=\bar\rho/(\mu m_p) && \text{cooling evaluated on coarse fields,}\\
\Ciso  &\equiv \left(\frac{\bar P}{k_B}\right)^{2}\int \mathcal{P}_V(T)\,\frac{\Leff(T)}{T^2}\,dT && \text{isobaric estimator,}
\end{align}
where $\mathcal{P}_V$ is the volume-weighted temperature PDF of the fine cells in the block. In practice $\Ciso$ is evaluated on a histogram with volume fractions $p_k$ and bin weights $w_k$:
\begin{equation}
\Cbin = \left(\frac{\bar P}{k_B}\right)^{2}\sum_k p_k\,w_k .
\end{equation}

\section{Exact relations}

\paragraph{Kernel identity.} For an ideal gas $n=P/(k_BT)$ exactly, so with
\begin{equation}
g(T)\equiv \frac{\Leff(T)}{T^2},\qquad n^2\Leff(T)=\frac{P^2}{k_B^2}\,g(T).
\end{equation}
This is the same kernel as in the emissivity PDF of Sharma et al.\ (Eq.~7), $\mathcal P_E\propto\mathcal P_V\Lambda/T^2$.

\paragraph{Conservation.} Because $\B$ is linear, $\sum_{ij}\Ctrue\,\ell^2=\int n^2\Leff\,dA=\dot\Sigma_{\rm cool}L_x$ exactly. Block averaging the cooling field loses nothing; all error lives in whatever model replaces $\Ctrue$.

\paragraph{Pressure fluctuations.} Write $P=\bar P(1+\pi)$ with $\B[\pi]=0$, and define the emissivity-weighted block average $\langle q\rangle_E\equiv\B[g\,q]/\B[g]$. Then, with no approximation,
\begin{equation}
\boxed{\;\frac{\Ctrue}{\Ciso}=1+2\langle\pi\rangle_E+\langle\pi^2\rangle_E\;}
\label{eq:isoerr}
\end{equation}
The leading correction is first order in $\pi$: it is the pressure fluctuation of the \emph{cooling} gas relative to the block mean.

\paragraph{Relation between $\Ciso$ and $\Ccg$.} Since $(\bar P/k_B)^2=\bar n^2\Tcg^2$,
\begin{equation}
\Ccg=\left(\frac{\bar P}{k_B}\right)^{2} g(\Tcg),\qquad
\frac{\Ciso}{\Ccg}=\frac{\B[g(T)]}{g(\Tcg)} .
\end{equation}
Both use the same pressure prefactor; $\Ccg$ evaluates the kernel at one temperature, $\Ciso$ averages it over the block's temperature distribution.

\paragraph{Exact form for $\Ccg$.} With $x\equiv1/T$ and $\tilde g(x)\equiv x^2\Leff(1/x)$, one has $\Tcg^{-1}=\B[(1+\pi)x]$ and
\begin{equation}
\boxed{\;\mathcal R\equiv\frac{\Ctrue}{\Ccg}=\frac{\B\!\left[(1+\pi)^2\,\tilde g(x)\right]}{\tilde g\!\left(\B[(1+\pi)x]\right)}\;}
\end{equation}
For isobaric blocks $\mathcal R=\B[\tilde g(x)]/\tilde g(\B[x])$, and by Jensen's inequality $\Ccg$ under-cools where $\tilde g$ is convex and over-cools where it is concave.

\section{Accuracy of $\Ccg$}

\subsection{Perturbative (smooth, single-phase) blocks}
Expand about $\hat x=\B[x]$ with $x=\hat x(1+\delta)$. Let $\alpha=d\ln\Lambda/d\ln T$, so $d\ln\tilde g/d\ln x=2-\alpha$, and define the curvature
\begin{equation}
K\equiv\frac{\hat x^2\tilde g''}{\tilde g}=(2-\alpha)(1-\alpha)+\frac{d\alpha}{d\ln T}.
\end{equation}
To second order,
\begin{equation}
\mathcal R\approx1+\frac{K}{2}\,\sigma_\delta^2+(2-\alpha)\,\B[\pi\delta]+\sigma_\pi^2,
\qquad \sigma_\delta^2=\B[\delta^2],\ \sigma_\pi^2=\B[\pi^2].
\end{equation}
In the isobaric limit $\delta=\delta n/\bar n$ and $\mathcal R\approx1+\tfrac K2\sigma_n^2$: the density clumping factor weighted by the curvature of the cooling curve. The expansion is invalid whenever the block's temperature range straddles a cutoff ($g$ is discontinuous there) or the block PDF is bimodal.

\subsection{Two-phase blocks}
Consider an isobaric block with cold volume fraction $f_c$ at $T_c$, hot gas at $T_h=\chi T_c$, and no intermediate gas. Then $\Ctrue=0$ exactly (both phases have cooling off), but
\begin{equation}
\Tcg=\frac{T_c}{f_c+(1-f_c)/\chi}
\end{equation}
falls inside the cooling window for
\begin{equation}
\frac{1/(0.95\chi)-1/\chi}{1-1/\chi}<f_c<\frac{1/1.05-1/\chi}{1-1/\chi},
\end{equation}
which for $\chi=100$ is $5.3\times10^{-4}<f_c<0.952$. Essentially every block containing both phases is assigned spurious cooling ($\mathcal R=0$), typically near the $\sim10^5$~K peak since the harmonic mean is pulled toward $T_c$. For comparison, $\B[T]$ would enter the window only for $f_c>0.05$ and would land near $0.9\,T_h$.

\subsection{Laminar $\tanh$ limit}
Suppose within a block $T$ varies only with height along the mean profile of Sharma et al.\ (Eq.~6), $\langle T\rangle=\tfrac{T_h-T_c}{2}\tanh(z/z_0)+\tfrac{T_h+T_c}{2}$. With $u=z/z_0$ the harmonic mean integrates in closed form:
\begin{equation}
\Tcg^{-1}=\frac{z_0}{T_c\,\ell}\left[u-\frac{\chi-1}{2\chi}\ln\!\left(1+\chi e^{2u}\right)\right]_{u_-}^{u_+}.
\end{equation}
The true row-averaged cooling follows from Eq.~8 of Sharma et al.,
\begin{equation}
\langle\Ctrue\rangle_j=\frac{\dot\Sigma_{\rm cool}}{\ell}\Big[\overline C_E(\langle T\rangle_+)-\overline C_E(\langle T\rangle_-)\Big],
\end{equation}
with $\overline C_E$ the cumulative emissivity distribution in mean temperature. For $\ell\ll z_0$, $\sigma_\delta^2\approx(\ell^2/12)(d\ln T/dz)^2$. For $\chi=100$, $d\ln T/dz$ peaks at $1.64/z_0$ at $T\approx10\,T_c$, right at the cooling peak, so $\sigma_\delta^2\lesssim0.22\,(\ell/z_0)^2$. In this limit $\Ccg$ is perturbatively correct (tens of percent) for $\ell\lesssim z_0$; for $\ell\gtrsim z_0$ the TRML occupies one or two rows and the two-phase result dominates. Horizontal (fractal) structure only adds variance, so this is a lower bound on the error of $\Ccg$.

\section{Accuracy of $\Ciso$}

$\Ciso$ removes the entire Jensen/two-phase problem of $\Ccg$: it is exact for any temperature distribution, including bimodal ones and ones that straddle the cutoffs, provided the block is isobaric. Its error has exactly two sources: pressure fluctuations within the block (Eq.~\ref{eq:isoerr}) and discretisation of the PDF.

\subsection{Isobaric error}

The error is $2\langle\pi\rangle_E+\langle\pi^2\rangle_E$. Only pressure variations \emph{within} a block matter; a compression that is uniform across the block is absorbed into $\bar P$ and causes no error. Three regimes arise.

\paragraph{Interface blocks.} Cooling gas in the TRML sits in the pressure dip supported by the compressive stress $\mathcal R_{zz}>0$ and is collapsing quasi-isobarically, but lags pressure equilibrium. If this gas is underpressured relative to the block mean, $\langle\pi\rangle_E<0$ and $\Ciso$ \emph{overestimates} the cooling. The size of the effect is set by the depth of the dip across one block and by $\mathcal M_{\rm turb}^2$; it should be a systematic bias at the level of $2|\langle\pi\rangle_E|$, largest in rows spanning the dip.

\paragraph{Acoustically triggered blocks.} Sound waves launched by collapsing clumps move stable-phase gas across the cutoffs, which is why cooling can switch on in a cell for one timestep and off the next. Here the cooling gas is, by construction, the perturbed gas, and the size of $\pi$ is fixed by the cutoffs. For an adiabatic wave ($\gamma=5/3$), $\delta\ln T=\tfrac{\gamma-1}{\gamma}\,\delta\ln P=0.4\,\delta\ln P$:
\begin{itemize}
\item Cold gas crossing $1.05\,T_c$ needs $\delta\ln P\ge0.122$, so $(1+\pi)^2\gtrsim1.28$ and $\Ciso$ \emph{underestimates} by $\gtrsim25\%$.
\item Hot gas crossing $0.95\,T_h$ needs $\delta\ln P\le-0.128$, so $(1+\pi)^2\lesssim0.77$ and $\Ciso$ \emph{overestimates} by $\gtrsim30\%$.
\end{itemize}
These are lower bounds on the local error, but they apply only when the wavelength is comparable to or smaller than $\ell$. For waves much longer than a block, $\pi\approx0$ and $\Ciso$ is exact. These blocks usually contribute little to the total cooling, but they dominate the \emph{time variability} of cooling in the stable phases.

\paragraph{Smooth single-phase blocks.} Here $\pi\sim\mathcal M^2$ and the correction is small.

\paragraph{Diagnostic.} The isobaric assumption can be tested per block without computing $\Ctrue$: with a volume PDF, Eq.~(\ref{eq:harmonic}) predicts $\sum_kp_k/T_k=\Tcg^{-1}$, while in general $\Tcg^{-1}=\B[(1+\pi)x]$. The fractional mismatch equals $\B[\pi x]/\B[x]$, a volume-and-$1/T$-weighted pressure covariance. This is a proxy for, but not identical to, the emissivity-weighted $\langle\pi\rangle_E$.

\paragraph{External evidence.} Sharma et al.\ find that $\mathcal P_E\propto\mathcal P_V\Lambda/T^2$, which is the same isobaric approximation applied across the whole TRML, agrees closely with the measured emissivity PDF (their Fig.~3). This indicates the isobaric error is small at the TRML-integrated level. It does not rule out larger per-block errors that cancel on average.

\subsection{Binning error}

Write $\B[g]=\sum_kp_k\langle g\rangle_k$, where $\langle g\rangle_k$ is the volume average of $g$ over the gas actually in bin $k$. Then the binning error is
\begin{equation}
\frac{\Cbin-\Ciso}{(\bar P/k_B)^2}=\sum_k p_k\left(w_k-\langle g\rangle_k\right).
\end{equation}

\paragraph{Smooth bins.} Within a bin of width $h=0.1\,{\rm dex}=0.230$ in $\ln T$, write $g\propto T^{\beta}$ with $\beta=\alpha-2$. If gas is uniform in $\ln T$ across the bin, a centre-point weight $w_k=g(T_k)$ underestimates $\langle g\rangle_k$ by $\beta^2h^2/24$:

\begin{center}
\begin{tabular}{ccc}
\toprule
$|\beta|$ & uniform-in-bin error & worst case (gas at one edge), $\sim|\beta|h/2$\\
\midrule
2 & 0.9\% & 23\%\\
4 & 3.5\% & 46\%\\
8 & 14\% & 92\%\\
\bottomrule
\end{tabular}
\end{center}

Large $|\beta|$ occurs on the steep low-temperature side of the cooling curve just above $10^4$~K. That is also where $g$ peaks (where $\alpha=2$) and where the emissivity PDF $\mathcal P_E$ peaks ($\sim10^{4.1}$~K in Sharma et al.). Bin-integrated weights, $w_k=h^{-1}\int_{\rm bin}g\,d\ln T$ computed from the cooling table, remove the second-order term. The first-order term, from non-uniform gas within a bin, remains and can only be reduced with narrower bins.

\paragraph{Cutoff bins (dominant problem).} With 40 log bins over $10^3$--$10^7$~K, the edges fall at multiples of $0.1$~dex, but the cutoffs do not:

\begin{center}
\begin{tabular}{lccc}
\toprule
cutoff & $\log_{10}T$ & containing bin & bin centre\\
\midrule
$1.05\,T_c$ & 4.021 & $[4.0,4.1]$ & $10^{4.05}=1.12\,T_c$ (cooling on)\\
$0.95\,T_h$ & 5.978 & $[5.9,6.0]$ & $10^{5.95}=0.89\,T_h$ (cooling on)\\
\bottomrule
\end{tabular}
\end{center}

All gas in $[T_c,1.05\,T_c]$ and $[0.95\,T_h,T_h]$ is therefore assigned cooling that the simulation does not apply. The cold side is critical. Any compressive acoustic fluctuation lifts cold gas above $T_c$ into this bin, which can put up to roughly half the cold-phase volume in the compressive half of a wave there. That gas is $\chi$ times denser than hot gas and is weighted at close to the peak of $g$. Integrated over a cold layer several parsecs deep, this spurious term can be comparable to or larger than the true $\dot\Sigma_{\rm cool}$, depending on how tightly the cold phase sits at $T_c$.

The outer bins ($10^3$--$10^4$ and $10^6$--$10^7$~K) are not empty: acoustic rarefactions and compressions populate them. They must be kept, with $w_k=0$. The fix is to place bin edges exactly at $1.05\,T_c$ and $0.95\,T_h$, so that every bin lies entirely inside or entirely outside the cooling window. Bins outside the window then carry exactly zero weight. Acoustically driven crossings into the window are counted correctly, apart from the isobaric error above.

\section{Summary}

\begin{center}
\renewcommand{\arraystretch}{1.3}
\begin{tabular}{>{\raggedright}p{2.4cm}>{\raggedright}p{5.4cm}>{\raggedright\arraybackslash}p{6.2cm}}
\toprule
Estimate & Exact relation to $\Ctrue$ & Accuracy\\
\midrule
$\Ccg$ & $\Ctrue/\Ccg=\B[(1+\pi)^2\tilde g(x)]/\tilde g(\B[(1+\pi)x])$ &
Perturbative $1+\tfrac K2\sigma_\delta^2+\dots$ only in smooth single-phase blocks; spurious cooling ($\mathcal R=0$) in essentially all two-phase blocks. Poor wherever $\ell\gtrsim z_0$.\\
$\Ciso$ (exact PDF) & $\Ctrue/\Ciso=1+2\langle\pi\rangle_E+\langle\pi^2\rangle_E$ &
Exact for isobaric blocks of any PDF shape. Systematic bias $2\langle\pi\rangle_E$ in interface blocks; $\gtrsim25$--$30\%$ in acoustically triggered blocks with wavelength $\lesssim\ell$.\\
$\Cbin$ & $\Ciso+(\bar P/k_B)^2\sum_kp_k(w_k-\langle g\rangle_k)$ &
Few percent with bin-integrated weights in smooth bins; tens of percent near $10^4$~K for spiky PDFs; $O(1)$ or worse with the current cutoff-straddling bins.\\
\bottomrule
\end{tabular}
\end{center}

\paragraph{How good is $\Ciso$?} Conceptually it is the correct estimator. It is exact up to within-block pressure fluctuations of the cooling gas, and it fixes the structural failure of $\Ccg$ at two-phase interfaces. Its accuracy is limited in practice by, in order of importance: (i) cutoff aliasing in the current binning, which is fixable and potentially dominant; (ii) the first-order isobaric bias $2\langle\pi\rangle_E$ in interface and acoustic blocks, which is irreducible for any isobaric method and has opposite signs in the two regimes; and (iii) within-bin variation of $g$ near $10^4$~K, which is reducible with bin-integrated weights and finer bins there.

\paragraph{Measurements to make on the fine data.} For each block and each $\Delta$, compute $\Ctrue$, $\Ciso$ (unbinned, from fine cells), and $\Cbin$ for the current and cutoff-aligned binnings, together with $\Ccg$. The quantities to plot are:
\begin{itemize}
\item $\ln(\Ciso/\Ctrue)$ against $\langle\pi\rangle_E$, $\Tcg$, height $z$ and cold fraction $f_c$. This is the isobaric floor, and the regime dependence identifies interface versus acoustic blocks.
\item $\ln(\Cbin/\Ciso)$ for both binnings. This is the discretisation floor, and the per-bin contribution isolates cutoff leakage.
\item $\ln(\Ccg/\Ctrue)$ as the baseline.
\item Totals $\sum\mathcal C\,\ell^2$ against $\dot\Sigma_{\rm cool}$, and row profiles $\langle\mathcal C\rangle_j(z)$, as functions of time. Acoustically triggered blocks should show up as the time-variable component.
\item The harmonic-mean mismatch $\sum_kp_k/T_k-\Tcg^{-1}$ as a cheap isobaricity check.
\end{itemize}

\paragraph{Reference.} P.~Sharma, A.~Kumar, D.~Datta, A.~Babul, R.~Das, K.~Aditya, \emph{Universal Structure of Turbulent Radiative Mixing Layers}, arXiv:2509.03802 (2025).

\end{document}